% Working manuscript: dense Strain-oriented reframing of the CMO study.
% The historical 9-page manuscript.tex is intentionally kept unchanged while
% this version is developed and audited.
\documentclass[11pt]{article}

\usepackage[T1]{fontenc}
\usepackage{lmodern}
\usepackage[margin=0.85in]{geometry}
\usepackage{graphicx}
\usepackage{booktabs}
\usepackage{amsmath,amssymb}
\usepackage{siunitx}
\usepackage{float}
\usepackage{microtype}
\usepackage{url}
\usepackage[colorlinks=true,linkcolor=blue!55!black,citecolor=blue!55!black,urlcolor=blue!55!black]{hyperref}
\usepackage[labelfont=bf,font=small]{caption}
\usepackage{subcaption}
\usepackage{enumitem}
\setlength{\emergencystretch}{2em}
\setlist{nosep}

\newcommand{\um}{\ensuremath{\mu\mathrm{m}}}
\newcommand{\ue}{\ensuremath{\mu\varepsilon}}
\newcommand{\RMS}{\mathrm{RMS}}
\newcommand{\SE}{\mathrm{SE}(3)}
\newcommand{\MaybeFigure}[2]{%
  \IfFileExists{#1}{\includegraphics[width=#2\linewidth]{#1}}{%
    \fbox{\parbox[c][0.19\textheight][c]{#2\linewidth}{\centering Figure generated by analysis workflow\\\texttt{#1}}}}}

\title{From ray-space calibration to strain-space validation:\\
scale-dependent error in a common-main-objective stereo microscope}
\author{Jean-Fran\c{c}ois Witz\\
\small Univ. Lille, CNRS, Centrale Lille, UMR 9013 -- LaMcube, France}
\date{Working draft -- \today}

\begin{document}
\maketitle

\begin{abstract}
Stereo digital image correlation ultimately measures spatial differences of reconstructed motion, yet stereo calibrations are commonly selected from image residuals or pointwise three-dimensional errors. This mismatch is particularly important for common-main-objective (CMO) stereo microscopes, whose non-central optical geometry also violates the single-centre assumption of conventional camera models. We combine a compact ray-space identification of the CMO geometry with an explicit strain-space validation on archived rigid translations. A flexible Zernike rayfield is first used as a geometric measuring instrument; residual structure then identifies a quasi-telecentric CMO model with per-channel rigid arm alignment. The resulting 26-parameter physical model reaches \SI{1.06}{px} reprojection RMS on the real Pycaso configuration, compared with more than \SI{300}{px} for the tested conventional stereo calibration. The mechanical validation uses two independent 101-plane translations and reserves 91 planes from each acquisition for evaluation. On the fine \SI{0.08}{mm} sweep, a direct stereo-to-object polynomial gives the smallest nominal-depth discrepancy (\SI{5.20}{\micro\metre}) but develops a persistent apparent-strain floor of approximately \SIrange{288}{301}{\micro\strain}, whereas the ray representation has a larger depth discrepancy but a strain floor of only \SIrange{50}{56}{\micro\strain}. The short-range endpoint-increment scales remain comparable (approximately \SIrange{0.20}{0.27}{\micro\metre}), showing that the difference is spatial structure rather than simply smaller point noise. Across gauge lengths and directions, the exact virtual-gauge strain is reproduced by the longitudinal increment of the rigid-alignment residual with correlations close to unity. This establishes the relevant calibration quantity for strain as the spatial increment of reconstruction error over the mechanical gauge, not the pointwise error itself. Rigid-target sequences can therefore test zero-strain consistency without requiring the rigid translation amplitude to be independently calibrated. The results support task- and scale-dependent validation of stereo calibrations before deformation measurement.
\end{abstract}

\textbf{Keywords:} stereo digital image correlation; stereo microscopy; calibration; apparent strain; ray field; non-central camera; virtual gauge; rigid-body motion.

% ======================================================================
\section{Introduction}
% ======================================================================

Stereo digital image correlation (stereo-DIC) maps image correspondences to a three-dimensional surface and then differentiates reconstructed motion to obtain strain. Calibration is upstream of both operations. It is therefore tempting to judge a calibration by the quantity that is easiest to compute---typically a pixel reprojection residual, a reconstructed target-coordinate residual, or a pointwise depth error. Those quantities are necessary diagnostics, but they are not the final mechanical observable. A strain estimate depends on differences between reconstruction errors at separated material points and on how those differences evolve between images. A spatially common reconstruction error can be large in position and nearly disappear from a length change; a smaller but spatially structured error can do the opposite.

The distinction is familiar in uncertainty propagation but is rarely made operational when a calibration model is selected. Reu \cite{reu2013} propagated calibration uncertainty through three-dimensional DIC and showed that calibration error must be considered at the level of reconstructed motion. The present work takes a complementary experimental route: instead of assigning a probability distribution to calibration parameters, we use rigid target motions for which the true strain is known exactly to be zero. The amplitude of the rigid translation need not itself be known to establish that condition. This permits the calibration chain to be interrogated directly in strain space.

Microscope stereo vision makes the problem especially visible. Schreier et al. \cite{schreier2004} established light-microscope stereo vision for deformation measurements, while polynomial mappings such as the Soloff formulation \cite{soloff1997} provide a successful black-box alternative and are implemented for microscope calibration in PYCASO \cite{pycaso2023}. At the same time, general imaging systems need not be central. Vision-ray calibration \cite{bothe2010} and modern generic camera models \cite{schops2020} represent an image position by a viewing ray rather than forcing all rays through a single optical centre. A common-main-objective stereo microscope is a concrete non-central case: the two channels share a main objective and view the object through off-axis sub-apertures. A conventional pinhole-plus-distortion parameterisation can therefore fail structurally rather than merely because it lacks enough distortion coefficients.

The CMO configuration considered here exposes two different questions that should not be conflated. The first is optical: how can a compact and interpretable calibration model be identified when direct joint estimation of optics and board poses is poorly conditioned? The second is mechanical: once several calibrations reconstruct the same rigid target, which one is most appropriate for deformation measurement? The first question motivates a flexible rayfield as an intermediate measuring instrument. The second requires observables that depend on spatial error increments, gauge length and direction.

This article therefore compresses the optical identification to the elements required for a mechanically credible calibration and develops the strain-space validation in detail. The principal contributions are as follows. First, a measured non-central rayfield is used to identify a compact quasi-telecentric CMO model with per-channel \SE{} corrections; the residual field is used to determine the missing physical degrees of freedom rather than only to report a fit error. Second, rigid translations are converted into a scale- and direction-dependent strain-consistency test. Third, the measured apparent-strain curves are linked quantitatively to the longitudinal structure of the non-rigid reconstruction residual. Finally, the internal calibration evidence is connected to an independent specimen/profilometry check, while explicitly separating coordinate conventions, rigid gauge freedoms and true metric differences.

The resulting claim is deliberately narrower than ``the ray model is more accurate.'' No representation is assumed to be universally superior. Instead, the calibration is evaluated in the coordinate and spatial scale of the intended mechanical measurement. The central proposition tested here is that \emph{pointwise calibration accuracy is neither necessary nor sufficient for strain accuracy; the relevant quantity is the spatial increment of reconstruction error over the mechanical gauge.}

% ======================================================================
\section{Ray-space identification of the CMO calibration}
\label{sec:cmo}
% ======================================================================

\subsection{Why the microscope is treated as non-central}

The instrument is a Leica M205C common-main-objective stereomicroscope from the archived PYCASO configuration. The two camera channels share the main objective but use distinct effective sub-pupils. In a central camera, every chief ray can be represented by a single centre $C$ and a direction field $d(u,v)$. In the CMO geometry, the more appropriate primitive is a line
\begin{equation}
\mathcal L_c(u,v)=\left\{O_c(u,v)+\lambda d_c(u,v):\lambda\in\mathbb R\right\},
\label{eq:rayline}
\end{equation}
where $c\in\{L,R\}$ denotes the channel. The ray origin need not collapse to one camera centre.

\begin{figure}[!htbp]
\centering
\MaybeFigure{../cmo/figures/cmo_physical.pdf}{0.80}
\caption{Common-main-objective stereo geometry used in the physical identification. Two off-axis effective sub-pupils share the main objective. The figure is inherited from the CMO analysis but only the effective chief-ray geometry is used here; the model is not presented as a full lens prescription.}
\label{fig:cmo_physical}
\end{figure}

On the tested real configuration, the standard central stereo calibration is not rescued by increasing conventional distortion complexity: the resulting reprojection error remains above \SI{300}{px}. This failure is used only as evidence that the central family is structurally inappropriate for this configuration, not as a claim that conventional calibration generally fails on microscopes.

\subsection{Measure rays first, identify optics second}

Directly fitting optical parameters and board poses creates a compensation problem: changes in the physical model can be partly absorbed by changes in the target poses. To separate these roles, the calibration proceeds in two stages. A flexible rayfield is first fitted to ChArUco observations. It is then treated as a geometric observable against which compact physical hypotheses are tested. This decomposition is the part of the CMO methodology retained in the present paper; detailed order sweeps, regularisation studies and alternative-family ablations are relegated to the reproducibility material.

The flexible field writes the origin and direction as low-order Zernike expansions over the image domain. For this data set the physical-construction reference uses a rigid origin field and a second-order direction field. The ray representation is gauge-fixed by selecting one point along each line; this removes a line-coordinate redundancy but does not turn the fitted rayfield into a metrological ground truth. Its role is to reveal geometric structure.

A candidate physical model is compared with the measured rayfield using direction and Pl\"ucker-moment residuals. The progression is diagnostic. A perspective CMO skeleton leaves a large structured angular residual. A quasi-telecentric skeleton removes most of it but leaves a low-order, nearly spatially uniform component. Allowing an independent rigid transform for each optical arm removes that component and yields the final 26-parameter CMO model. Thus the per-arm \SE{} terms are not added because they simply reduce a scalar residual; they are introduced because the residual topology identifies a global arm-alignment degree of freedom.

\begin{figure}[!htbp]
\centering
\MaybeFigure{../cmo/figures/residual_evolution.pdf}{0.96}
\caption{Compressed physical-identification sequence. The structured residual of the wrong perspective family motivates the quasi-telecentric description; the remaining low-order residual motivates per-channel rigid arm alignment; the final CMO+\SE{} model leaves a near-noise-floor ray-space residual. The full hypothesis and ablation study is retained in the repository rather than repeated in the main article.}
\label{fig:residual_evolution}
\end{figure}

The resulting physical model contains 14 quasi-telecentric parameters plus six rigid-arm parameters per channel, for 26 in total. On the ten CMO calibration views it gives a pixel RMS of \SI{1.06}{px} (median \SI{0.87}{px}, 95th percentile \SI{1.84}{px}), compared with \SI{0.47}{px} for the 57-parameter flexible rayfield. The lower flexible-model residual is expected and is not interpreted as evidence that the flexible reconstruction is metrically better. The compact model is retained because it captures the observed non-central geometry with fewer, physically interpretable degrees of freedom.

\subsection{Observability and the role of the flexible rayfield}

The optical parameters and board poses are not all independently observable from near-parallel planar views. A Schur-complement analysis of the information matrix shows that several optical directions become weak after pose marginalisation. This matters because a direct bundle adjustment can improve pixel RMS while moving along optics--pose compensation directions. The flexible rayfield partly avoids that failure mode by separating geometric measurement from physical interpretation. In the long-form CMO study this point is explored through several priors and bundle-adjustment variants; here it is retained only as a limitation on parameter interpretation. Effective baseline, working distance and convergence angle are meaningful descriptors within the selected rayfield gauge and reference scale, but they are not all independent absolute metrological measurements.

This observability issue will reappear in Section~\ref{sec:specimen}: the near-parallel calibration poses constrain lateral geometry much more strongly than axial scale. Crucially, a weakly observed axial metric is different from an arbitrary non-rigid shape disagreement.

% ======================================================================
\section{Mechanical validation framework}
\label{sec:mechanical}
% ======================================================================

\subsection{Rigid motion as a zero-strain experiment}

Let $\widehat{\mathbf X}_0(\mathbf x)$ be the reconstructed reference position associated with a target location $\mathbf x$, and let $\widehat{\mathbf X}_t(\mathbf x)$ be its reconstruction after a rigid motion. The true object satisfies
\begin{equation}
\mathbf X_t(\mathbf x)=\mathbf R_t\mathbf X_0(\mathbf x)+\mathbf T_t,
\qquad \mathbf R_t\in SO(3).
\end{equation}
The value of $\mathbf T_t$ need not be known to establish that every true material distance is invariant.

For each reconstructed frame we therefore fit the best proper rigid transform from the reconstructed reference cloud to the current cloud and define the non-rigid residual
\begin{equation}
\mathbf e_t(\mathbf x)=
\widehat{\mathbf X}_t(\mathbf x)-
\left[\mathbf R_t\widehat{\mathbf X}_0(\mathbf x)+\mathbf T_t\right].
\label{eq:rigidresidual}
\end{equation}
No scale and no reflection are allowed in this fit. Equation~\eqref{eq:rigidresidual} is therefore the deformation that the reconstruction chain invents after the rigid part has been removed. It should not be confused with the coordinate-convention reflection discussed later for the independent coin comparison: that fixed reflection relates two model coordinate systems and is removed before comparing their surfaces; it is not used to improve the rigid-target strain scores.

\subsection{Exact apparent strain of a virtual gauge}

Consider target points $i$ and $j$. The exact reconstructed gauge strain at frame $t$ is
\begin{equation}
\varepsilon^{\mathrm{app}}_{ij,t}=
\frac{\left\|\widehat{\mathbf X}_{j,t}-\widehat{\mathbf X}_{i,t}\right\|}
{\left\|\widehat{\mathbf X}_{j,0}-\widehat{\mathbf X}_{i,0}\right\|}-1.
\label{eq:exactstrain}
\end{equation}
Because the denominator is the reconstructed reference length, Eq.~\eqref{eq:exactstrain} tests length-change consistency rather than manufactured target scale. For a rigid object its exact expected value is zero, independently of the stage displacement.

The earlier short analysis used only gauges spanning four target pitches (nominally \SI{1.2}{mm}). Here all admissible separations are used in the horizontal, vertical and two diagonal directions. This exposes a property that a single RMS cannot show: calibration error can change character with gauge length.

\subsection{First-order strain from the spatial increment of reconstruction error}

Let
\begin{equation}
\mathbf q_{ij,0}=\widehat{\mathbf X}_{j,0}-\widehat{\mathbf X}_{i,0},
\qquad
L_{ij}=\|\mathbf q_{ij,0}\|,
\end{equation}
and let the rigidly transported unit gauge direction be
\begin{equation}
\mathbf t_{ij,t}=\frac{\mathbf R_t\mathbf q_{ij,0}}{L_{ij}}.
\end{equation}
Substituting Eq.~\eqref{eq:rigidresidual} into the current gauge vector gives
\begin{equation}
\widehat{\mathbf X}_{j,t}-\widehat{\mathbf X}_{i,t}
=
\mathbf R_t\mathbf q_{ij,0}
+
\left[\mathbf e_t(\mathbf x_j)-\mathbf e_t(\mathbf x_i)\right].
\end{equation}
Linearising its norm with respect to the residual increment yields
\begin{equation}
\delta\varepsilon_{ij,t}
\simeq
\frac{\mathbf t_{ij,t}^{T}}{L_{ij}}
\left[\mathbf e_t(\mathbf x_j)-\mathbf e_t(\mathbf x_i)\right].
\label{eq:firstorder}
\end{equation}
Equation~\eqref{eq:firstorder} is the mechanical core of the paper. Pointwise error $\|\mathbf e\|$ does not appear by itself. Only its longitudinal spatial increment appears. A common-mode offset cancels; a coherent spatial gradient survives as strain; short-range uncorrelated contributions are progressively averaged by longer gauges.

\subsection{A scale- and direction-dependent structure measure}

For gauges of length $L$ and orientation $\theta$, define the longitudinal second-order structure amplitude
\begin{equation}
D_{\parallel}(L,\theta)=
\left\langle
\left(
\mathbf t^T
[\mathbf e(\mathbf x+L\mathbf t)-\mathbf e(\mathbf x)]
\right)^2
\right\rangle,
\label{eq:structure}
\end{equation}
where the average is over admissible gauge locations and held-out rigid motions. Equation~\eqref{eq:firstorder} then gives
\begin{equation}
\sigma_\varepsilon(L,\theta)
\simeq
\frac{\sqrt{D_{\parallel}(L,\theta)}}{L}.
\label{eq:structurestrain}
\end{equation}
This formulation supplies a direct bridge between a spatial calibration residual and a mechanical strain floor.

For compact interpretation we also fit the empirical two-term law
\begin{equation}
\sigma_\varepsilon^2(L)=
\left(\frac{A}{L}\right)^2+B^2.
\label{eq:scalelaw}
\end{equation}
Parameter $A$ has dimensions of displacement and captures an endpoint-increment contribution that decays approximately as $1/L$. Parameter $B$ is a gauge-length-independent coherent strain floor. Equation~\eqref{eq:scalelaw} is not used to fit the calibration; it is a post-hoc diagnostic of the measured rigid-motion curves.

% ======================================================================
\section{Data, calibration representations and evaluation protocol}
\label{sec:protocol}
% ======================================================================

\subsection{Archived CMO microscope translations}

The PYCASO archive provides two series of 101 stereo image pairs from the same CMO microscope configuration family \cite{pycasodata}. Images are $2048\times2048$ pixels and contain a ChArUco target with $16\times12$ squares of nominal side \SI{0.3}{mm}, hence 165 internal corners. The wide series spans nominal depths from \SI{2.65}{mm} to \SI{3.35}{mm} in \SI{7}{\micro\metre} increments; the fine series spans \SIrange{2.96}{3.04}{mm} in \SI{0.8}{\micro\metre} increments.

The target pose and image scale are visibly different between the two acquisitions. They are therefore treated as separate calibration experiments rather than as training and test sequences of an unchanged camera state. For the wide series, ten planes with zero-based indices $0,10,20,30,40,50,60,80,90,100$ are used for calibration. Ten approximately uniformly spaced planes are used for the fine series. The remaining 91 planes of each acquisition are reserved for evaluation.

Corners are detected directly with OpenCV ChArUco detection and subpixel refinement. No missing corner is synthetically completed for the three retrospective comparison estimators. This yields 13,444 held-out wide-series corner observations and 15,015 held-out fine-series observations. The compact CMO identification uses the denoised/completed observation preparation developed in the physical-model study on the ten wide calibration frames; its held-out mechanical test below is nevertheless performed on the same raw reserved correspondences. This front-end difference is reported explicitly because it prevents treating every model row as an identical estimator fit.

\subsection{Comparison representations}

Four roles are distinguished.

The \emph{finite-plane ray field} represents, independently for each camera, the intersection of a viewing ray with a plane $z$ by
\begin{equation}
\begin{pmatrix}X\\Y\end{pmatrix}
=\mathbf a_c(u,v)+\zeta\,\mathbf b_c(u,v),
\qquad
\zeta=\frac{z-z_0}{s_z},
\label{eq:finiteplane}
\end{equation}
with complete real Zernike bases for the components of $\mathbf a_c$ and $\mathbf b_c$. Reconstruction is the midpoint of the closest points on the two rays.

The \emph{forward Soloff-type model} fits image coordinates as polynomial functions of the object coordinates and reconstructs a point by numerical inversion. The \emph{direct model} fits object coordinates directly as a polynomial function of the four stereo image coordinates. The direct representation does not define independent camera rays but provides an important control: it can be extremely accurate for one reconstructed coordinate without necessarily preserving distances.

Finally, the \emph{CMO 26p model} is the compact physical model of Section~\ref{sec:cmo}. It is evaluated as the physically interpreted outcome of the rayfield-mediated workflow, not as another member of the degree-selected polynomial benchmark.

\subsection{Model selection and held-out observables}

For the three retrospective comparison estimators, polynomial/Zernike order is selected by leaving out one of the ten calibration planes at a time and minimising plane-balanced depth RMS. The selected representation is refitted on all ten calibration planes and never sees the remaining 91 evaluation planes during selection.

Three families of held-out observables are then reported. First, pointwise depth discrepancy relative to the nominal file label measures consistency with the assumed stage coordinate. Second, reference-based centroid motion measures how reconstructed rigid translation compares with the nominal depth change. Third, Eqs.~\eqref{eq:exactstrain}--\eqref{eq:scalelaw} measure rigid-body strain consistency without requiring the stage displacement to be exact. The third family is the principal mechanical validation.

A separate spatial holdout, using alternating target corners, is retained as a supplementary interpolation check. Synthetic distorted-perspective controls are also retained to verify implementation correctness and to demonstrate that the ray representation does not intrinsically suppress imposed extension.

% ======================================================================
\section{Results}
\label{sec:results}
% ======================================================================

\subsection{Pointwise depth and strain select different models}

The first result reproduces the ranking reversal that motivated the mechanical analysis. On the wide \SI{0.70}{mm} sweep, the ray and forward reconstructions have nearly identical held-out depth discrepancies, \SI{11.61}{\micro\metre} and \SI{11.59}{\micro\metre}, while the direct model gives \SI{12.36}{\micro\metre}. At a nominal \SI{1.2}{mm} gauge, the three apparent-strain levels are of the same order and no general ray advantage is observed.

The fine sweep is qualitatively different. The direct affine mapping gives the smallest pointwise depth discrepancy, \SI{5.20}{\micro\metre}; the Soloff-type model gives \SI{5.89}{\micro\metre}; and the ray field gives \SI{9.41}{\micro\metre}. Yet the combined horizontal/vertical \SI{1.2}{mm} apparent-strain RMS is approximately \SI{184}{\micro\strain} for the ray field, \SI{268}{\micro\strain} for the forward model and \SI{338}{\micro\strain} for the direct mapping. Thus the calibration that best follows nominal fine-sweep depth is the worst of the three for preserving rigid-target distances.

That observation alone would establish only a data-set-specific ranking reversal. The following analyses identify the spatial mechanism and its dependence on the mechanical gauge.

\subsection{Gauge length exposes a coherent strain floor}

\begin{figure}[!htbp]
\centering
\MaybeFigure{figures/gauge_length_dependence.pdf}{0.96}
\caption{Apparent rigid-motion strain RMS versus virtual-gauge length for the wide and fine acquisitions. Gauge length is varied over every admissible horizontal and vertical target separation; the plotted axial result combines the two orientations only for compact visual comparison. The fine sweep separates the models strongly as gauge length increases.}
\label{fig:gaugelength}
\end{figure}

On the fine sweep, short gauges initially give similar error magnitudes. For horizontal gauges of \SI{0.3}{mm}, the ray, Soloff and direct models give respectively 765, 894 and 753~\ue. Increasing the gauge to \SI{2.4}{mm} changes these values to 108, 187 and 311~\ue. Vertical gauges show the same behaviour: at \SI{2.4}{mm} the corresponding values are 99, 166 and 301~\ue. The direct mapping therefore stops benefiting from gauge enlargement and approaches a persistent false-strain level near \SI{300}{\micro\strain}; the ray representation continues to decrease approximately as $1/L$ over most of the available range.

The two-term decomposition in Eq.~\eqref{eq:scalelaw} quantifies this difference. On the fine sweep, the short-range displacement scale $A$ is similar for all three models: \SI{0.229}{\micro\metre} (ray), \SI{0.265}{\micro\metre} (Soloff) and \SI{0.207}{\micro\metre} (direct) horizontally; the vertical values are \SI{0.207}{\micro\metre}, \SI{0.261}{\micro\metre} and \SI{0.198}{\micro\metre}. The coherent strain floor $B$, however, differs strongly. It is 56/50~\ue{} (horizontal/vertical) for the ray field, 141/151~\ue{} for Soloff and 301/288~\ue{} for the direct mapping. The fits have $R^2>0.998$ in squared-RMS space.

\begin{table}[!htbp]
\centering
\caption{Fine-sweep gauge-scale decomposition, $\sigma_\varepsilon^2=(A/L)^2+B^2$. Parameter $A$ is a short-range longitudinal endpoint-increment scale; $B$ is the coherent apparent-strain floor.}
\label{tab:scalelaw}
\begin{tabular}{lrrrr}
\toprule
& $A_H$ (\um) & $B_H$ (\ue) & $A_V$ (\um) & $B_V$ (\ue)\\
\midrule
Ray field & 0.229 & 55.8 & 0.207 & 49.6\\
Soloff-type & 0.265 & 140.7 & 0.261 & 150.9\\
Direct polynomial & 0.207 & 300.6 & 0.198 & 288.0\\
\bottomrule
\end{tabular}
\end{table}

This decomposition changes the interpretation of the original ranking result. The ray field is not simply winning because every reconstructed point is quieter. At the shortest scale its endpoint-increment magnitude is comparable to the alternatives. The decisive difference is the amount of spatially coherent distortion that survives length averaging.

The wide sweep does not show the same separation. Its fitted coherent floors are closer and the preferred model depends on direction. This is important: the observed scale behaviour is not an algebraic consequence of using rays. It is a property of the error field produced by a given calibration on a given acquisition. That is precisely why task-specific validation is required.

\subsection{The spatial-increment equation predicts the measured strain}

\begin{figure}[!htbp]
\centering
\MaybeFigure{figures/structure_prediction.pdf}{0.96}
\caption{Exact virtual-gauge strain RMS and its first-order prediction from the longitudinal spatial increment of the rigid-alignment residual. The near overlap is an experimental verification of Eq.~\eqref{eq:firstorder}, not a fitted calibration relation.}
\label{fig:structureprediction}
\end{figure}

The mechanism predicted by Eq.~\eqref{eq:firstorder} is directly testable because both sides can be computed independently after reconstruction. On the fine sweep at \SI{1.2}{mm}, the correlation between exact gauge strain and the first-order spatial-increment prediction is 0.99998 for the ray field, 0.99937 for the Soloff-type model and effectively 1.00000 for the direct mapping in the horizontal direction. The corresponding vertical and diagonal correlations are likewise close to unity, and the agreement improves further at longer gauges.

The ranking reversal can therefore be attributed to measured spatial error structure rather than to an abstract distinction between model families. The direct mapping has excellent nominal-depth behaviour but its residual field contains a coherent longitudinal component. Because this component grows approximately in proportion to gauge separation, division by $L$ does not remove it and the strain tends to a non-zero floor. The ray-field residual contains a larger fraction of short-range increment error that is reduced by increasing $L$.

\subsection{Spatial maps make the false strain visible}

\begin{figure}[!htbp]
\centering
\MaybeFigure{figures/pseudo_strain_maps.pdf}{0.96}
\caption{Representative fine-sweep apparent-strain maps for nominally \SI{1.2}{mm} horizontal and vertical gauges. The target undergoes rigid translation, so every non-zero value is reconstruction-induced pseudo-strain. A common colour scale is used across model columns.}
\label{fig:pseudostrainmaps}
\end{figure}

The fine-sweep maps show why a scalar pointwise RMS is insufficient. The ray-field pseudo-strain varies spatially around zero with relatively little large-scale bias. The forward polynomial retains a visible coherent component. The direct mapping produces a broad same-sign field over much of the target. Such a smooth field can coexist with low pointwise depth error because a spatial gradient is inexpensive in position while remaining directly visible to a length measurement.

This observation also motivates reporting direction. The stereo baseline and non-central optical geometry provide preferred directions, and the calibration residual need not be isotropic. Horizontal, vertical and diagonal gauge families are therefore retained separately in the machine-readable results even when two of them are combined in compact figures.

\subsection{Held-out mechanical behaviour of the compact CMO model}

The 26-parameter physical CMO model was identified on the ten wide-series calibration frames. It is consequently evaluated mechanically only on the 91 reserved planes of that acquisition; it is not transferred to the fine acquisition, whose image pose and scale differ and which is independently calibrated throughout this paper. The same raw held-out correspondences are triangulated with the CMO rays and assessed with the rigid-residual procedure above.

% Numerical paragraph is filled automatically/manually from
% results/cmo_mechanical_observables.json after the workflow audit.
\textbf{CMO held-out result placeholder.} The analysis workflow writes the full gauge-length curves, ray-gap statistics and $(A,B)$ decomposition to \texttt{results/cmo\_mechanical\_observables.json}. This paragraph will be replaced by audited numbers before submission; no claim about the CMO strain ranking is made from its calibration reprojection RMS alone.

\begin{figure}[!htbp]
\centering
\MaybeFigure{figures/cmo_wide_gauge_length.pdf}{0.94}
\caption{Held-out wide-sweep apparent strain versus gauge length including the compact CMO 26p physical calibration. The CMO model is evaluated only on the acquisition from which it was identified.}
\label{fig:cmo_holdout}
\end{figure}

\subsection{Synthetic controls: low pseudo-strain is not strain suppression}

A calibration representation could appear artificially favourable on rigid motion if its reconstruction suppressed all changes in length. Controlled distorted-perspective simulations exclude this explanation. With no coordinate noise, all tested models recover a prescribed \SI{1000}{\micro\strain} horizontal extension with an RMS strain error below \SI{0.03}{\micro\strain}. Under coordinate noise, the models exhibit comparable strain noise rather than a generic ray-field advantage. The real-data scale separation is therefore not a built-in consequence of representing the cameras by rays.

% ======================================================================
\section{Independent surface check and coordinate/gauge separation}
\label{sec:specimen}
% ======================================================================

The rigid-target analysis measures consistency but cannot determine the absolute axial scale of a reconstruction. The CMO study therefore includes an independent speckled coin specimen and profilometric relief of the same embossed region. This check is retained because it addresses a different question from the rigid-motion strain test.

The CMO and Zernike surface reconstructions are first compared only after accounting for coordinate convention and rigid gauge. The CMO implementation uses the documented $v\mapsto Y$ sign convention, which produces a fixed $Y$-axis reflection relative to the comparison frame. A reflection has determinant $-1$ and cannot be absorbed by an \SE{} fit. It is therefore corrected explicitly first. A subsequent proper rigid alignment removes an approximately $8.6^\circ$ pose difference. After these operations the median dense three-dimensional point-to-point discrepancy falls from approximately \SI{3.5}{mm} to \SI{0.018}{mm}, and the plane-normalised depth residual is approximately \SI{0.007}{mm}. The gross pre-alignment disagreement is thus a convention/pose effect, not evidence of arbitrary non-rigid surface mismatch.

\begin{figure}[!htbp]
\centering
\MaybeFigure{../cmo/figures/zernike_cmo_rigid_removed.pdf}{0.96}
\caption{Separation of coordinate convention, rigid pose and residual metric difference in the independent specimen reconstruction. The $Y$ reflection is a coordinate convention and is removed before proper rigid alignment. The residual difference is predominantly axial scale, not an arbitrary non-rigid deformation.}
\label{fig:rigidremoved}
\end{figure}

After convention and rigid effects are removed, the remaining relief amplitudes differ by a depth-only factor: lateral scales satisfy approximately $s_x\simeq s_y\simeq1$, whereas $s_z\simeq0.70$, consistent with a Zernike-to-CMO relief ratio near 1.4. The near-parallel calibration poses leave this axial degree of freedom comparatively weak. Stereo data alone therefore do not establish which amplitude is physically correct.

External profilometry resolves that ambiguity for the measured coin region. The profilometric relief standard deviation is \SI{20.8}{\micro\metre}; the Soloff reconstruction gives \SI{22.8}{\micro\metre}, the CMO 26p reconstruction \SI{23.0}{\micro\metre}, and the flexible Zernike reconstruction \SI{34.0}{\micro\metre}. All retain a shape correlation of about 0.96 to the profilometric relief. On this specimen the CMO and Soloff amplitudes are therefore within roughly 10\% of the external reference, while the flexible rayfield over-amplifies relief. This does not convert the coin into a traceable full-volume calibration artifact, but it demonstrates why low ray-space residual alone must not be equated with metric reconstruction accuracy.

The specimen and rigid-target experiments are complementary. Profilometry constrains the axial metric for one independent surface; rigid translation constrains the strain consistency of the reconstruction field over multiple positions and gauge lengths. Neither should be substituted for the other.

% ======================================================================
\section{Discussion}
% ======================================================================

\subsection{What quantity should a DIC user validate?}

The experimental results make Eq.~\eqref{eq:firstorder} operational. If the intended observable is strain, the relevant calibration statistic is not simply $\RMS(\|\mathbf e\|)$ or a pixel reprojection RMS. It is the longitudinal increment of $\mathbf e$ over the spatial support of the strain measurement. This has a direct practical implication: a calibration report should include a rigid-motion apparent-strain curve versus gauge length, and preferably versus direction, at scales comparable to the DIC virtual strain gauge.

This conclusion does not make pointwise metrics irrelevant. Depth and centroid-motion discrepancies remain essential when out-of-plane displacement is the observable. The fine-sweep direct mapping illustrates the distinction particularly well: it is superior for nominal depth and reference-based centroid motion while inferior for long-gauge zero-strain consistency. The appropriate ranking therefore depends on what the experiment is intended to measure.

The scale-law decomposition adds a useful diagnostic vocabulary. A large $A$ with small $B$ indicates short-range endpoint noise that can be reduced by increasing gauge length. A large $B$ indicates a coherent distortion field that survives gauge enlargement and is therefore more dangerous for strain. Two calibrations with similar local displacement noise can have very different strain floors if their spatial covariance differs.

\subsection{Why rigid translations are especially useful}

An independently calibrated translation stage would be required to interpret nominal depth discrepancy as absolute 3D accuracy. It is not required to establish zero true strain of a rigid target. This asymmetry is valuable for laboratory validation: a sequence of rigid target images can reveal reconstruction-induced strain even when the exact rigid motion is uncertain. The approach is therefore experimentally lighter than a full strain standard, while being more directly relevant to deformation measurement than reprojection RMS.

There are limits. Target rigidity and thermal stability must still be credible. The present ChArUco target dimensions and flatness are assumed rather than traceably certified, and corner localisation is not identical to dense speckle matching. A rigid-target strain-consistency test should consequently be read as a property of the complete fiducial/reconstruction chain under the tested imaging conditions, not as a complete uncertainty budget for a future loaded-speckle experiment.

\subsection{What the rayfield contributes---and what it does not}

The flexible rayfield plays two distinct roles. In the optical identification it is a measuring instrument that exposes the structure missing from a compact physical hypothesis. In the mechanical comparison it is one possible reconstruction representation. These roles must not be merged. Its low calibration residual does not make it a metric ground truth; the coin profilometry explicitly shows a case in which the flexible reconstruction has the wrong axial gain despite tight ray intersections.

Conversely, the compact CMO model should not be declared mechanically superior merely because it is physically interpretable or because it reaches \SI{1.06}{px}. Its mechanical credibility comes from held-out rigid-motion and independent-surface checks. This distinction is one reason to combine the previously separate CMO and strain arguments in a single, shorter optical narrative but a richer mechanical validation.

\subsection{Generalisability}

The particular 26-parameter CMO skeleton is instrument-specific, but the validation framework is not. Any stereo reconstruction method that returns matched three-dimensional points can be subjected to Eqs.~\eqref{eq:rigidresidual}--\eqref{eq:structurestrain}. The same analysis can compare central camera models, generic rayfields, polynomial mappings and multi-camera systems without forcing their fitting residuals into a common likelihood space.

The current evidence covers two acquisitions of one microscope, not a population of CMO instruments. The wide and fine sequences also differ in image pose and scale, so their different gauge-scale behaviour should not be attributed uniquely to translation amplitude. Further validation should vary magnification, target pose and instrument, and should include a loaded speckled specimen with an independent strain reference. Those additions would test transferability; they are not required for the present demonstration that pointwise and differential calibration errors are distinct observables.

% ======================================================================
\section{Conclusion}
% ======================================================================

A stereo calibration intended for deformation measurement should be validated in strain space, not selected from reprojection or pointwise depth error alone. On the CMO microscope studied here, a rayfield-mediated procedure identifies a compact non-central physical model and provides the geometric basis for reconstruction, while held-out rigid translations reveal the mechanical consequences of calibration error.

The decisive result is spatial. On the fine translation series, three reconstruction models have comparable short-range endpoint-increment scales of order \SI{0.2}{\micro\metre}, yet their coherent apparent-strain floors differ by about a factor of six: approximately \SIrange{50}{56}{\micro\strain} for the ray field, \SIrange{141}{151}{\micro\strain} for the Soloff-type model and \SIrange{288}{301}{\micro\strain} for the direct mapping. The direct mapping is nevertheless the best of the three in nominal depth. Exact gauge strain is predicted almost perfectly by the longitudinal spatial increment of the rigid-alignment residual, confirming the mechanism rather than merely observing a ranking reversal.

Rigid-target sequences consequently provide a practical zero-strain consistency test whose interpretation does not require exact knowledge of the imposed rigid displacement. The appropriate calibration metric is scale- and direction-dependent and should be matched to the intended mechanical gauge. Independent profilometry remains necessary when absolute surface scale, rather than differential strain consistency, is the quantity of interest.

\section*{Data and code availability}
The archived PYCASO images are referenced at their original revision \cite{pycasodata}. All analysis scripts, saved observations, prediction arrays, generated mechanical-observable tables and figure-generation code are maintained in the StereoComplex repository \cite{analysiscode}. The branch used for the present reframing preserves the earlier short manuscript and adds independent scripts for gauge-length analysis, rigid-residual structure, scale-law decomposition and held-out CMO evaluation.

\bibliographystyle{unsrt}
\bibliography{references}

\end{document}
